\section*{Within the Citadel}\stepcounter{section}\phantomsection\addcontentsline{toc}{section}{Within the Citadel}
\begin{DndReadAloud}
	Beyond the gates, the expedition enters a vast interior hall of dark stone and tarnished metal. Broad stairways rise toward elevated galleries while narrow bridges cross the open spaces between them. Most of the citadel remains without power, illuminated only by the occasional glow of a surviving conduit or rune.

	Several chambers branch from the main route, their doors standing partially open after centuries of neglect. Farther within, faint traces of energy continue deeper into the citadel.
\end{DndReadAloud}
{\noindent\entryfont The expedition can explore the chambers branching from the central hall before continuing deeper into the citadel. Prototype Assembly and Defence Cabinet are optional locations and may be omitted if time is running short. Neither is required to reach the Core Room, but both provide additional insight into Aberdour and its final years.}
\subsection*{\inTextCircled[10]{11}{10} Prototype Assembly}\stepcounter{subsection}\phantomsection\addcontentsline{toc}{subsection}{Prototype Assembly}
{\entryfont A heavy metal door branches from the central hall, its surface marked with faded workshop sigils and several dormant warning runes. Unlike most of the citadel, a weak magical seal still clings to the entrance.

\paragraph*{Opening the Door} Normally, the door is held shut by a weak magical seal. It can be opened with a successful \textbf{DC 14 Intelligence (Arcana) Check} and \textbf{DC 14 Dexterity (Sleight of Hand) Check} using Thieves' Tools to disrupt or bypass the remaining enchantment.

If an \textbf{Emergency Discharge} occurred while opening the citadel, the magical seal has already collapsed, but the ancient locking mechanism has seized completely. The door can then be forced open with a \textbf{DC 14 Strength (Athletics) Check}, repaired or bypassed with an appropriate Tinker's Tools check, or opened effortlessly with L.A.T.C.H's Ghostkey.}
\begin{DndReadAloud}
	Beyond the door lies a broad workshop crowded with stone benches, suspended frames, and the unfinished remains of ancient constructs.

Some are little more than articulated limbs and exposed mechanisms. Others resemble crude ancestors of the machines found throughout Auchtertool, their designs growing steadily more sophisticated from one workbench to the next.

Notes, diagrams, and engraved plates cover almost every available surface.
\end{DndReadAloud}
{\noindent\entryfont The surviving records concern experimental construct frames, power regulators, mobility systems, and numerous projects that were apparently never completed. Most are of historical rather than immediate practical value, but they reveal how rapidly construct design advanced within Aberdour.

On one stone desk lies a collection of notes written in Dwarvish. They record the unexplained disappearance of several precision tools, prototype components, and important blueprints. Repeated annotations suggest the loss was never resolved before the citadel was abandoned.}

\subsection*{\inTextCircled[10]{12}{10} Defence Cabinet}\stepcounter{subsection}\phantomsection\addcontentsline{toc}{subsection}{Defence Cabinet}
{\entryfont A massive reinforced door seals one of the chambers branching from the central hall. Unlike the magically secured Prototype Assembly, no enchantment remains upon it - the problem is simply its weight. Centuries of neglect have left the hinges and locking bars almost completely seized.

\paragraph*{Opening the Door} A character can force the door open with a successful DC 17 Strength (Athletics) check. Suitable tools, leverage, or other creative methods may grant advantage or provide another means of entry.}
\begin{DndReadAloud}
	Beyond the heavy door lies a sombre strategic chamber. A broad stone table dominates its centre, surrounded by cabinets, map racks, and shelves filled with brittle records.

	Most of the documents have been reduced to blackened scraps. Those that survive depict the Mines of Methven and the City of Aberdeen.

	Whatever happened here, someone took considerable effort to ensure that much of this room's knowledge did not survive.
\end{DndReadAloud}
\subsubsection*{Military Records}
{\entryfont Most of the strategic records have been deliberately burned. Whether this occurred during Aberdour's decline, as an act of sabotage, or as an attempt to keep sensitive information from falling into the wrong hands cannot be determined.

One partially preserved collection concerns a project referred to only as an \textbf{"Ultimate Weapon of War"}. The records describe attempts to develop an entirely new form of power supply - an experimental energy core capable of sustaining something vastly more demanding than an ordinary construct.}